% Source: https://artofproblemsolving.com/community/c3h1564751_geometry_problem
% Stable identity: aops-topic-1564751-q01
Put $A=(0,0)$, $B=(2,0)$, $C=(2,2)$, $D=(0,2)$. Since $EB=ED$, $E$ lies on the perpendicular bisector of $BD$, the line $y=x$; since $C$ is inside triangle $BDE$, $E=(t,t)$ with $t>2$.

\textbf{The point $F$.} Line $DE$ through $(0,2)$ and $(t,t)$ meets the line $x=2$ (line $BC$) at
\[
F=\left(2,\ 2+\dfrac{2(t-2)}{t}\right)=\left(2,\ 4-\dfrac{4}{t}\right),
\]
above $C$: $F$ is on side $DE$ and on line $BC$ extended beyond $C$.

\textbf{The parallel condition.} $FA\parallel EB$ means the vectors $F-A=\left(2,4-\dfrac{4}{t}\right)$ and $E-B=(t-2,t)$ are parallel:
\[
2t=\left(4-\dfrac{4}{t}\right)(t-2)\iff2t^2=(4t-4)(t-2)\iff t^2-6t+4=0 .
\]
So $t=3\pm\sqrt5$, and $t>2$ forces $t=3+\sqrt5$. Then $\dfrac{4}{t}=\dfrac{4}{3+\sqrt5}=3-\sqrt5$ and $F=\left(2,1+\sqrt5\right)$.

\begin{asy}
usepackage("amsmath");
size(6cm);
defaultpen(fontsize(10pt));
real t=3+sqrt(5);
pair A=(0,0),B=(2,0),C=(2,2),D=(0,2),Ep=(t,t),F=(2,1+sqrt(5));
draw(A--B--C--D--cycle,linewidth(0.9));
draw(B--D--Ep--cycle,linewidth(0.9));
draw(C--F,red+linewidth(1));
draw(A--F,blue+dashed);
draw(B--Ep,blue);
draw(B--(C+0*(1,0)),linewidth(0.9));
dot(F);dot(C);
label("$A$",A,SW);label("$B$",B,SE);label("$C$",C,(1.2,-0.6));label("$D$",D,W);label("$E$",Ep,NE);label("$F$",F,NW);
label("$2$",(A+B)/2,S);label("$2$",(A+D)/2,W);
pair m1=(B+Ep)/2,m2=(D+Ep)/2;
draw(m1+0.15*unit(rotate(90)*(Ep-B))--m1-0.15*unit(rotate(90)*(Ep-B)));
draw(m2+0.15*unit(rotate(90)*(Ep-D))--m2-0.15*unit(rotate(90)*(Ep-D)));
label("\mbox{}",(min(currentpicture,user=true)+max(currentpicture,user=true))/2,black); // keeps the text after the figure black when LaTeX typesets the labels inline
\end{asy}

Hence
\[
FC=\left(1+\sqrt5\right)-2=\boxed{\sqrt5-1}\text{ cm}\approx1.236\text{ cm}.
\]
(Check: $FA$ and $EB$ both have slope $\dfrac{1+\sqrt5}{2}$.)
